\section{Conclusion}
\label{sec:conclusion}

We ran a paired fault-injection study of a LangGraph ReAct agent on WebArena under
the official evaluator, and its most useful product is a negative result about the
measurement rather than about the agent.

End-to-end success is a blunt, high-variance instrument. With
\ResWebDomMissingNNat--\ResWebTimeoutNNat{} native pairs per fault and an observed
discordance of \PiDNative{}, no minimum detectable effect is defined at \Power{}
power, no fault is significant after Holm correction, and detecting a $10$~pp
difference would need on the order of \PairsForTenPPObserved{} pairs per fault. The
same faults are unambiguous at the process level --- HTTP errors alone add
\BehWebHttpErrorStepDelta{} steps and \BehWebHttpErrorTimeDelta{}~s per pair ---
which is where robustness measurement should look. Layered on top of that is a
warning about instruments: a compatibility fallback failed by construction on all
\OfficialFallbackN{} trials it handled and shifted the frozen baseline's headline
rate by \BaselineSwing{}~pp, more than any effect we set out to measure, while the
recovery artifact consisted mostly of constant columns and an outcome restated as a
behavioural label. Neither would have been visible from a single success number.

The model-versus-architecture question remains open: answering it requires a
crossed design under the same paired protocol, and we release the staged plan for it
alongside this paper. What we can say now is narrower and, we think, more durable:
for web agents, report the paired design, stratify the evaluator, state the power,
and measure the process --- because the success rate alone will not tell you whether
your agent is robust or merely unmeasured.
